\subsection{Contraste funcional independiente de los núcleos}
El modelo de estimación delimita los treinta componentes de los núcleos Operacional, Servicios y Administrativo. Un caso representa un objetivo observable del actor; las transacciones son intercambios completos de solicitud y respuesta. Las validaciones internas, clics y sentencias de persistencia no reciben puntos propios. Salida y regreso, preparación y participación de clase, consulta de salud y retorno, espera y cambio contractual se cuentan por objetivos distintos. Las alternativas de rechazo, duplicación o desconexión pertenecen al mismo objetivo. El modelo conserva actor, flujo, fuente y correspondencia con los noventa paquetes iniciales en \path{componentes/ucp_casos_nucleos.csv}.

La frontera incluye reglas y vistas propias de 1.4--1.6. Excluye construcción de identidad, canales, adaptadores y sincronización comunes, infraestructura, migración histórica, pilotos, formación y operación; estas partidas conservan su imputación ascendente. Los sistemas contables y dispositivos quedan fuera de la frontera de este conteo: sus adaptadores se estiman en 1.7. Los roles humanos se cuentan una sola vez aunque utilicen varios canales. El modelo es una descomposición de diseño de la oferta, no un levantamiento validado ni productividad observada.

La clasificación de casos y actores y los ajustes técnicos y ambientales siguen el método UCP descrito por Cohn; sus factores se aplican al modelo particular de la marina.~\parencite{cohn-ucp} Se aplica el método de puntos de casos de uso: actores simples/medios/complejos con pesos 1/2/3; casos con hasta tres, cuatro a siete y más de siete transacciones con pesos 5/10/15. $UUCP=UAW+UUCW$, $TCF=0,6+0,01\sum w_iT_i$ y $EF=1,4-0,03\sum w_iE_i$. Los trece factores técnicos y ocho de ambiente tienen valoración y justificación individual en \path{componentes/ucp_factores_nucleos.csv}. Las capacidades de ambiente son condiciones de provisión propuestas; no acreditan experiencia de personas ni contratos ya ejecutados.

El conteo contiene 34 casos y 12 roles complejos: $UAW=36$, $UUCW=180$ y $UUCP=216$. Se obtiene $TCF=1,165$, $EF=1,01$ y $UCP=254,16$. Con 1 factor de ambiente desfavorable se usa $CF=20$ HH/punto como referencia no calibrada.

Se declara la interpretación de esfuerzo de programación: $E_P=UCP\,CF$, y el esfuerzo total de referencia es $E_T=E_P/0,40$. Como supuesto de planificación de Synaptix, la distribución de referencia es análisis 10\,\%, diseño 20\,\%, programación 40\,\%, pruebas 15\,\% y sobrecarga 15\,\%. El reparto 10/20/40/15/15 y la lectura de programación son decisiones explícitas de esta estimación; no se atribuyen a una productividad histórica de Synaptix ni a una distribución universal del método UCP. Es una comparación independiente; no se agrega ese total al de la EDT ni se vuelven a sumar gestión y pruebas ya imputadas. Para evitar duplicar actores entre etapas, su peso se distribuye proporcionalmente al peso funcional de E1 y E2; los dos subtotales conservan exactamente el conteo global.

La Tabla~\ref{tab:sd7-ucp-contraste} presenta contraste de esfuerzo directo y referencia UCP.
\begin{table}[H]
\centering
\begin{tabularx}{\linewidth}{Y R{29mm} R{29mm}}\hline
 \EncabezadoTabla{Magnitud} & \EncabezadoTabla{EDT directa} & \EncabezadoTabla{Referencia UCP}\\\hline
Programación/Desarrollo de los núcleos & 2.400 HH & 5.083,13 HH\\\hline
Esfuerzo de los noventa paquetes / total de referencia & 4.200 HH & 12.707,82 HH\\\hline
\end{tabularx}
\caption{Contraste de esfuerzo directo y referencia UCP}
\label{tab:sd7-ucp-contraste}
\FuenteTabla{Elaboración propia: conteo de actores y casos, factores técnicos y ambientales y EDT inicial de los núcleos.}
\end{table}\par\smallskip
El contraste de programación alcanza 2,12 veces las HH de Desarrollo iniciales. La comparación de totales requiere además imputar análisis, diseño, pruebas y gobierno compartidos del mismo alcance; no se atribuye toda esa diferencia a subestimación porque la EDT directa excluye trabajos comunes. Aun así, la diferencia de programación no queda explicada por la mera existencia de bibliotecas reutilizables: no se dispone de un factor de productividad medido que permita reducir $CF$ hasta coincidir con la EDT.

La sensibilidad deteriora E1 y E3 de 3 a 2 y E7 de 3 a 4, manteniendo las restantes valoraciones. Con cuatro factores desfavorables, $EF=1,115$ y $CF=28$, la programación asciende a 7.856,20 HH y el total de referencia a 19.640,50 HH. Los registros de contraste por etapa y sensibilidad permiten reproducir estos resultados.

\paragraph{Modelo de interacción y regla de conteo.} El conteo se basa en el modelo de interacción propuesto para la oferta. Cada transacción identifica una entrada del actor y la respuesta que devuelve el sistema antes de esperar otro estímulo. Validar permisos, comprobar restricciones, calcular resultados y persistir datos son procesamiento de esa respuesta, salvo que se identifique otra entrada independiente. Una confirmación sólo se cuenta cuando el flujo requiere una decisión posterior del actor; no se añade una petición de validación para elevar el peso. El detalle de los 34 objetivos se conserva en \path{componentes/ucp_transacciones_nucleos.csv} y en su desarrollo detallado del Formulario T-15. Son decisiones de diseño propuestas, sujetas a validación con los usuarios, no interacciones ya observadas. Los rechazos dentro de una respuesta no reciben puntos adicionales; las resoluciones que exigen otra entrada están identificadas en el flujo.
\ifdefined\EnCuerpoSD
El detalle de las 34 transacciones por objetivo se presenta en el Formulario T-15, Tabla~\textit{Transacciones del modelo funcional propuesto}; T-14 conserva su correspondencia con los paquetes. En el capítulo, el conteo y las tablas de esfuerzo permiten seguir el cálculo sin repetir el listado completo.
\else
La Tabla~\ref{tab:t15-ucp-transacciones} documenta las entradas y respuestas que sustentan el peso de cada objetivo.
\begingroup\small\begin{longtable}{L{14mm}L{27mm}L{97mm}R{9mm}}
\caption{Transacciones del modelo funcional propuesto}\label{tab:t15-ucp-transacciones}\\\hline
 \EncabezadoTabla{Caso} & \EncabezadoTabla{Objetivo} & \EncabezadoTabla{Entradas y respuestas independientes} & \EncabezadoTabla{Peso}\\\hline\endfirsthead
\hline \EncabezadoTabla{Caso} & \EncabezadoTabla{Objetivo} & \EncabezadoTabla{Entradas y respuestas independientes} & \EncabezadoTabla{Peso}\\\hline\endhead
CU-01 & Registrar salida & 1. Actor: Identifica embarcación y patrón.\newline Sistema: Presenta contexto de salida y formulario.\newline 2. Actor: Presenta nómina y consentimiento y solicita registrar salida.\newline Sistema: Valida condición de registro y devuelve salida y hora registradas. & 5\\\hline
CU-02 & Confirmar regreso seguro & 1. Actor: Identifica nave e informante.\newline Sistema: Presenta salida pendiente.\newline 2. Actor: Presenta evidencia de regreso.\newline Sistema: Contrasta reporte y presenta discrepancias.\newline 3. Actor: Resuelve discrepancia y confirma regreso.\newline Sistema: Registra estado y responsable. & 5\\\hline
CU-03 & Declarar plan voluntario & 1. Actor: Identifica nave y responsable.\newline Sistema: Presenta formulario del plan.\newline 2. Actor: Presenta destino, duración y consentimiento.\newline Sistema: Valida y devuelve plan registrado. & 5\\\hline
CU-04 & Atender plan vencido & 1. Actor: Consulta vencimiento.\newline Sistema: Presenta plan vencido y evidencia.\newline 2. Actor: Registra contacto o intento fallido.\newline Sistema: Devuelve contacto registrado.\newline 3. Actor: Decide escalamiento según regla.\newline Sistema: Registra actuación y estado. & 5\\\hline
CU-05 & Autorizar amarra & 1. Actor: Consulta nave y puesto.\newline Sistema: Presenta compatibilidad y ocupación.\newline 2. Actor: Solicita asignación.\newline Sistema: Comprueba exclusividad local y confirma o rechaza.\newline 3. Actor: Presenta evidencia de maniobra segura.\newline Sistema: Devuelve evidencia registrada. & 5\\\hline
CU-06 & Preparar clase y nómina & 1. Actor: Selecciona clase y monitor.\newline Sistema: Presenta ficha de clase.\newline 2. Actor: Incorpora alumnos y habilitaciones.\newline Sistema: Comprueba consentimientos y restricciones y devuelve nómina revisable.\newline 3. Actor: Publica nómina.\newline Sistema: Devuelve nómina publicada. & 5\\\hline
CU-07 & Registrar participación por intervalo & 1. Actor: Consulta clase y nómina.\newline Sistema: Presenta participantes.\newline 2. Actor: Registra alumno, bote, monitor e intervalo.\newline Sistema: Valida solapamientos y devuelve cambios registrados. & 5\\\hline
CU-08 & Confirmar retiro de alumno & 1. Actor: Presenta solicitud y autorizado.\newline Sistema: Comprueba consentimiento e identidad y devuelve resultado.\newline 2. Actor: Coordina punto de entrega.\newline Sistema: Presenta coordinación registrada.\newline 3. Actor: Confirma entrega.\newline Sistema: Devuelve entrega única registrada. & 5\\\hline
CU-09 & Consultar salud mínima de participantes & 1. Actor: Identifica monitor y clase y consulta salud.\newline Sistema: Comprueba participación efectiva y devuelve datos mínimos autorizados. & 5\\\hline
CU-10 & Conciliar retorno de clase & 1. Actor: Presenta retorno por bote.\newline Sistema: Contrasta participantes y presenta diferencias.\newline 2. Actor: Resuelve diferencia o registra alerta.\newline Sistema: Devuelve conciliación o alerta registrada.\newline 3. Actor: Confirma retorno completo.\newline Sistema: Devuelve cierre de retorno. & 5\\\hline
CU-11 & Actualizar condición de infraestructura & 1. Actor: Identifica puesto, boya o fondeo.\newline Sistema: Presenta estado vigente.\newline 2. Actor: Presenta condición y evidencia.\newline Sistema: Valida relaciones y vigencia y devuelve inventario actualizado. & 5\\\hline
CU-12 & Presentar inspección desconectada & 1. Actor: Identifica boya y orden.\newline Sistema: Presenta formulario de inspección.\newline 2. Actor: Registra hallazgos y evidencia sin enlace.\newline Sistema: Devuelve resultado guardado localmente.\newline 3. Actor: Solicita sincronización.\newline Sistema: Devuelve recepción o conflicto.\newline 4. Actor: Resuelve duplicado o conflicto.\newline Sistema: Devuelve conciliación registrada. & 10\\\hline
CU-13 & Asociar medidor temporalmente & 1. Actor: Identifica medidor e instalación.\newline Sistema: Presenta asociaciones vigentes.\newline 2. Actor: Presenta intervalo de asociación.\newline Sistema: Valida exclusividad temporal y devuelve asociación trazable. & 5\\\hline
CU-14 & Consultar habilitación inicial de faena & 1. Actor: Identifica faena o contratista y consulta habilitación.\newline Sistema: Comprueba vigencia y autorización y devuelve estado y bloqueos. & 5\\\hline
CU-15 & Programar faena de varadero & 1. Actor: Presenta nave, trabajo y fechas.\newline Sistema: Presenta puestos compatibles y restricciones de campaña y maniobra.\newline 2. Actor: Selecciona puesto y solicita reserva.\newline Sistema: Devuelve reserva o alternativa. & 5\\\hline
CU-16 & Autorizar izaje & 1. Actor: Identifica orden y nave.\newline Sistema: Presenta orden y requisitos de izaje.\newline 2. Actor: Presenta plan, medios y responsables.\newline Sistema: Comprueba restricciones y devuelve estado de revisión.\newline 3. Actor: Autoriza o suspende con evidencia.\newline Sistema: Devuelve decisión registrada. & 5\\\hline
CU-17 & Cerrar orden de trabajo & 1. Actor: Consulta orden.\newline Sistema: Presenta orden y habilitación.\newline 2. Actor: Presenta avance, incidencias y evidencia y solicita cierre.\newline Sistema: Valida resultado y devuelve cierre o necesidad de reprogramación. & 5\\\hline
CU-18 & Habilitar acceso a intervención & 1. Actor: Presenta identidad y orden.\newline Sistema: Comprueba acreditación, inducción, vigencia y zona y devuelve autorización o denegación.\newline 2. Actor: Registra entrada.\newline Sistema: Devuelve entrada registrada.\newline 3. Actor: Registra salida.\newline Sistema: Devuelve salida registrada. & 5\\\hline
CU-19 & Registrar despacho de combustible & 1. Actor: Identifica nave y operador.\newline Sistema: Presenta contexto del despacho.\newline 2. Actor: Presenta producto, volumen y evidencia en momento seguro.\newline Sistema: Valida y devuelve hecho local único. & 5\\\hline
CU-20 & Rectificar despacho con conciliación & 1. Actor: Consulta despacho.\newline Sistema: Presenta registro y evidencia.\newline 2. Actor: Presenta rectificación motivada.\newline Sistema: Comprueba autoridad y presenta diferencias.\newline 3. Actor: Confirma rectificación.\newline Sistema: Devuelve nueva versión conservando origen. & 5\\\hline
CU-21 & Conciliar serie de consumo & 1. Actor: Selecciona medidor e intervalo.\newline Sistema: Presenta lecturas, calidad y asociación temporal.\newline 2. Actor: Resuelve faltantes y anomalías.\newline Sistema: Devuelve serie conciliada.\newline 3. Actor: Publica serie.\newline Sistema: Devuelve serie validada publicada. & 5\\\hline
CU-22 & Consultar consumo e historia ambiental & 1. Actor: Solicita período y contrato.\newline Sistema: Valida identidad y pertenencia y devuelve serie común y cobertura. & 5\\\hline
CU-23 & Recibir residuo con trazabilidad & 1. Actor: Identifica origen y categoría.\newline Sistema: Presenta datos de recepción.\newline 2. Actor: Presenta cantidad, fecha y evidencia.\newline Sistema: Presenta discrepancias.\newline 3. Actor: Concilia discrepancias y confirma recepción.\newline Sistema: Devuelve recepción y retención registradas. & 5\\\hline
CU-24 & Cerrar entrega a gestor autorizado & 1. Actor: Selecciona residuos y gestor.\newline Sistema: Presenta saldo y autorización del gestor.\newline 2. Actor: Presenta entrega parcial.\newline Sistema: Valida cantidades y devuelve entrega registrada.\newline 3. Actor: Vincula certificado sanitario.\newline Sistema: Devuelve certificado vinculado.\newline 4. Actor: Confirma cierre.\newline Sistema: Devuelve saldo y expediente cerrado. & 10\\\hline
CU-25 & Publicar indicador e imputación de consumo & 1. Actor: Selecciona serie validada.\newline Sistema: Comprueba contrato y permanencia y devuelve indicador y cargo propuesto.\newline 2. Actor: Revisa exclusiones y resuelve diferencias.\newline Sistema: Devuelve propuesta ajustada.\newline 3. Actor: Publica resultado.\newline Sistema: Devuelve resultado trazable publicado. & 5\\\hline
CU-26 & Mantener maestro contractual & 1. Actor: Identifica persona, organización o nave.\newline Sistema: Presenta maestro vigente.\newline 2. Actor: Presenta datos y documentos.\newline Sistema: Valida duplicados y relaciones y devuelve versión y autoridad. & 5\\\hline
CU-27 & Consultar contrato vigente & 1. Actor: Identifica contrato y período y solicita consulta.\newline Sistema: Comprueba permisos y vigencia y devuelve documentos y límites. & 5\\\hline
CU-28 & Atender visita náutica & 1. Actor: Presenta nave visitante y responsable.\newline Sistema: Presenta disponibilidad comercial.\newline 2. Actor: Solicita autorización operacional.\newline Sistema: Devuelve autorización o rechazo de la autoridad local.\newline 3. Actor: Confirma asignación transitoria.\newline Sistema: Devuelve asignación y condiciones. & 5\\\hline
CU-29 & Gestionar lista de espera & 1. Actor: Presenta solicitud y prioridades.\newline Sistema: Valida elegibilidad y presenta posición o alternativa.\newline 2. Actor: Acepta alternativa o mantiene espera.\newline Sistema: Devuelve estado registrado. & 5\\\hline
CU-30 & Formalizar cambio contractual & 1. Actor: Identifica contrato y cambio.\newline Sistema: Presenta estado contractual.\newline 2. Actor: Presenta documentos y fechas y confirma cambio.\newline Sistema: Comprueba exclusividad operacional y devuelve versión contractual. & 5\\\hline
CU-31 & Presentar lote de hechos facturables & 1. Actor: Selecciona hechos conciliados.\newline Sistema: Aplica tarifa y comprueba duplicados y período y presenta lote.\newline 2. Actor: Presenta lote contable.\newline Sistema: Devuelve aceptación o rechazo registrado. & 5\\\hline
CU-32 & Conciliar respuesta contable & 1. Actor: Presenta emisión, pagos y saldo recibidos.\newline Sistema: Valida referencia y origen y presenta diferencias por moneda.\newline 2. Actor: Resuelve diferencia o solicita reintento.\newline Sistema: Devuelve conciliación o resultado del reintento.\newline 3. Actor: Confirma conciliación.\newline Sistema: Devuelve estado importado actualizado. & 5\\\hline
CU-33 & Resolver medida de cobranza & 1. Actor: Consulta deuda.\newline Sistema: Presenta deuda y evidencia.\newline 2. Actor: Presenta fundamento y propuesta.\newline Sistema: Devuelve propuesta pendiente de decisión.\newline 3. Actor: Registra decisión competente.\newline Sistema: Devuelve actuación sin sanción automática. & 5\\\hline
CU-34 & Exportar expediente de custodia & 1. Actor: Selecciona expediente y alcance y solicita exportación.\newline Sistema: Comprueba permisos y piezas, genera formatos y manifiesto, verifica integridad y PDF/A y devuelve exportación trazable. & 5\\\hline
\end{longtable}\endgroup

\FuenteTabla{Elaboración propia de Synaptix: modelo de interacción y reglas de conteo UCP.}
El peso se asigna por intercambios independientes, sin sumar validaciones internas; los registros preservan la correspondencia con componentes y paquetes del mismo alcance.
\fi

\paragraph{Esfuerzo seleccionado y conciliación.} Synaptix adopta la referencia UCP sin reducir el factor de conversión por reutilización no medida: 7.766 HH en E1 y 4.942 HH en E2, 12.708 HH en total. El redondeo por etapa deja la programación en 3.107 y 1.977 HH, respectivamente, por encima de los 3.106,36 y 1.976,77 HH de referencia. Las 4.200 HH originales se acreditan una sola vez y se añaden 8.508 HH en 216 paquetes de 8--80 HH, identificados en \path{componentes/paquetes_conciliacion_ucp.csv}. Cada lote tiene resultado, responsable y aceptación; sus precedencias completan especificación/diseño, construcción, verificación y documentación antes del manifiesto de versión.

La imputación por componente es proporcional a su peso funcional; el peso de los actores conserva el reparto E1/E2 del modelo. Análisis acredita las HH Funcional existentes; diseño acredita Datos, Integración y Seguridad propios del núcleo; programación acredita Desarrollo; pruebas acredita Calidad; sobrecarga cubre coordinación y documentación propias del componente. La matriz \path{componentes/conciliacion_ucp_actividades.csv} concilia exactamente el esfuerzo seleccionado. La seguridad compartida, los canales, adaptadores, migración, UAT integral, formación y operación permanecen fuera de esta frontera; las comprobaciones del núcleo no se presupuestan otra vez como UAT integral.

La Tabla~\ref{tab:sd7-ucp-seleccion} presenta distribución del esfuerzo seleccionado por actividad y etapa.
\begin{table}[H]
\centering
\begin{tabularx}{\linewidth}{Y R{21mm}R{21mm}R{22mm}}\hline
 \EncabezadoTabla{Actividad} & \EncabezadoTabla{E1 HH} & \EncabezadoTabla{E2 HH} & \EncabezadoTabla{Total HH}\\\hline
Análisis & 777 & 494 & 1.271\\\hline
Diseño & 1.553 & 988 & 2.541\\\hline
Programación & 3.107 & 1.977 & 5.084\\\hline
Pruebas & 1.165 & 741 & 1.906\\\hline
Sobrecarga & 1.164 & 742 & 1.906\\\hline
Total & 7.766 & 4.942 & 12.708\\\hline
\end{tabularx}
\caption{Distribución del esfuerzo seleccionado por actividad y etapa}
\label{tab:sd7-ucp-seleccion}
\FuenteTabla{Elaboración propia: conciliación UCP y redondeo por etapa, sin duplicar esfuerzo compartido.}
\end{table}\par\smallskip
La decisión fija el esfuerzo de diseño de la oferta y sus recursos recalculados; no acredita productividad histórica ni aprobación humana o contractual. La sensibilidad de 19.640,50 HH es un escenario de deterioro, separado del esfuerzo seleccionado: si cambian los factores de provisión, Proyecto recalculará esfuerzo, perfiles y calendario antes de comprometer los puestos afectados. El modelo delimita los núcleos, y no se presenta como conteo UCP de todo el software de las cuentas 1.3, 1.7 y 1.9.
